\subsection{Grundprinzip des Profilmodells}
\label{subsec:grundprinzip-profilmodell}

Wie in Kapitel~\ref{sec:architektur} beschrieben, bildet ein
Master-Repository die gemeinsame technische Ausgangsbasis. Das Modell weist
damit Merkmale featurebasierter Variabilitätsverwaltung und systematischer
Softwarewiederverwendung auf, ohne allein deshalb eine vollständige
Software-Produktlinie zu konstituieren
\parencite{clements2001software-product-lines,krueger1992software-reuse}.
Der Katalog \path{profiles/features.toml} trennt vier Begriffe: Der
\emph{Core} umfasst die in jedes Produktprojekt kopierten Pfade
\path{README.md}, \path{VERSION}, \path{.gitignore}, \path{config/},
\path{docs/}, \path{profiles/}, \path{shared/} und \path{tools/}. Ein
\emph{Feature} ordnet einer technischen Fähigkeit eigene Scaffold-Pfade und
gegebenenfalls Abhängigkeiten zu. Ein \emph{Profil} ist dagegen ein benanntes
Preset zulässiger Features. Eine \emph{optionale Capability} erweitert dieses
Preset nach der Profilwahl; sie ist weder ein Plattformprofil noch automatisch
Bestandteil eines Profils \parencite{kleiveist2026profiles,kleiveist2026template}.

\begin{figure}[H]
  \centering
  \resizebox{0.92\textwidth}{!}{%
  \begin{tikzpicture}[node distance=8mm]
    \node[casePrimary,minimum width=54mm] (master) {Master Template};
    \node[caseBox,minimum width=54mm,below=of master] (core) {Shared Core};
    \node[caseBox,minimum width=54mm,below=of core] (profiles) {Profile Presets};
    \node[caseOptional,minimum width=54mm,below=of profiles] (project) {Generated Product Project};

    \draw[caseFlow] (master) -- (core);
    \draw[caseFlow] (core) -- (profiles);
    \draw[caseFlow] (profiles) -- (project);
    \node[font=\sffamily\scriptsize,text=black!60,right=13mm of core,align=left]
      {gemeinsame Wartung\\keine separaten Template-Repositories};
  \end{tikzpicture}}
  \caption{Profilmodell auf Basis eines gemeinsamen Master-Templates}
  \label{fig:profile-model}
\end{figure}


Abbildung~\ref{fig:profile-model} fasst diese Ebenen zusammen. Sie abstrahiert
von den Laufzeitkomponenten und hebt stattdessen hervor, dass Core,
Feature-Katalog und Presets in einer Wartungsbasis liegen. Der Generator löst
das gewählte Preset und kompatible Capabilities auf, bildet aus Core- und
Feature-Pfaden einen Scaffold-Plan und kopiert diesen in ein neues
Produktprojekt. Nicht aktivierte Verzeichnisse werden somit überwiegend gar
nicht ausgewählt; bei Webprofilen entfernt die Nachbearbeitung zusätzlich die
Tauri-CLI aus den kopierten Frontend-Abhängigkeiten. Optional können Name,
Slug und bei Tauri die Anwendungs-ID angepasst werden
\parencite{kleiveist2026profiles,kleiveist2026tooling}.

Die Auswahl in Abbildung~\ref{fig:profile-selection-flow} berücksichtigt, dass
die gegenwärtigen Presets Backend und Cloud gemeinsam führen. Bei
Desktopprojekten trennt der zusätzliche Browser-Produktkanal semantisch
\texttt{full-platform} von \texttt{desktop-cloud}; erst nach der Profilwahl
wird eine Capability geprüft. Die anschließende Abhängigkeitsauflösung
verhindert, dass dadurch stillschweigend ein fehlendes Plattformfeature
ergänzt wird.

\begin{figure}[H]
  \centering
  \resizebox{\textwidth}{!}{%
  \begin{tikzpicture}[node distance=7mm and 8mm]
    \node[casePrimary,minimum width=35mm] (requirements) {Projektanforderung};
    \node[caseBox,minimum width=24mm,right=of requirements] (browser) {Browser?};
    \node[caseBox,minimum width=24mm,right=of browser] (desktop) {Desktop?};
    \node[caseBox,minimum width=24mm,right=of desktop] (backend) {Backend?};
    \node[caseBox,minimum width=24mm,right=of backend] (cloud) {Cloud?};
    \node[casePrimary,minimum width=35mm,below=13mm of $(desktop)!0.5!(backend)$] (profile) {passendes Profil};
    \node[caseOptional,minimum width=35mm,below=of profile] (capability) {optionale Capability};

    \draw[caseFlow] (requirements) -- (browser);
    \draw[caseFlow] (browser) -- (desktop);
    \draw[caseFlow] (desktop) -- (backend);
    \draw[caseFlow] (backend) -- (cloud);
    \draw[caseFlow] (browser) |- (profile);
    \draw[caseFlow] (desktop) -- (profile);
    \draw[caseFlow] (backend) -- (profile);
    \draw[caseFlow] (cloud) |- (profile);
    \draw[caseFlow] (profile) -- (capability);
  \end{tikzpicture}}
  \caption{Entscheidungsstruktur für die Auswahl eines Projektprofils}
  \label{fig:profile-selection-flow}
\end{figure}


Jedes generierte Projekt erhält das aktive Manifest
\texttt{project-\allowbreak profile.toml}. Es speichert ausschließlich
\texttt{schema\_version}, Profil-\texttt{id}, \texttt{name},
\texttt{description}, die angeforderten \texttt{optional\_features} und die
vollständig aufgelösten \texttt{features}. Das gemeinsame Tooling lädt und
validiert das Manifest. Dadurch behandelt es nur aktive Dienste und Workflows.
Im Master aktiviert das Manifest \texttt{full-platform} und
\texttt{postgres}. Das gleichnamige Preset bleibt davon unberührt.
Strukturprofil und Laufzeitumgebung bleiben getrennt: Development-, Test- und
Produktionswerte, Ports, URLs und Geheimnisse stammen aus dem in
Abschnitt~\ref{subsec:konfigurationsmodell} beschriebenen
Konfigurationsmodell \parencite{kleiveist2026configuration,kleiveist2026template}.

Abbildung~\ref{fig:profile-matrix} ermöglicht den schnellen strukturellen
Vergleich. Tabelle~\ref{tab:profile-matrix} ergänzt dazu die exakt deklarierten
Feature-IDs und die daraus ausgewählten Laufzeitverzeichnisse. Insbesondere ist
PostgreSQL bei keinem der fünf Presets voreingestellt.

\begin{figure}[H]
  \centering
  \resizebox{\textwidth}{!}{%
  \begin{tikzpicture}[
    cell/.style={draw=caseLine,minimum width=24mm,minimum height=9mm,align=center,font=\sffamily\scriptsize},
    head/.style={cell,fill=caseBlue,text=white,font=\sffamily\bfseries\scriptsize},
    active/.style={cell,fill=caseBlue!10,text=caseBlue,font=\sffamily\bfseries\scriptsize},
    optional/.style={cell,fill=caseGreen!12,text=caseGreen!55!black,font=\sffamily\bfseries\scriptsize}
  ]
    \matrix[matrix of nodes,nodes in empty cells,ampersand replacement=\&,row sep=-\pgflinewidth,column sep=-\pgflinewidth] {
      |[head,minimum width=32mm]| Profil \& |[head]| Frontend \& |[head]| FastAPI \& |[head]| Tauri \& |[head]| Cloud \& |[head]| PostgreSQL \\
      |[head,minimum width=32mm]| web-only \& |[active]| enthalten \& |[cell]| -- \& |[cell]| -- \& |[cell]| -- \& |[cell]| -- \\
      |[head,minimum width=32mm]| web-cloud \& |[active]| enthalten \& |[active]| enthalten \& |[cell]| -- \& |[active]| enthalten \& |[optional]| optional \\
      |[head,minimum width=32mm]| desktop-local \& |[active]| enthalten \& |[cell]| -- \& |[active]| enthalten \& |[cell]| -- \& |[cell]| -- \\
      |[head,minimum width=32mm]| desktop-cloud \& |[active]| enthalten \& |[active]| enthalten \& |[active]| enthalten \& |[active]| enthalten \& |[optional]| optional \\
      |[head,minimum width=32mm]| full-platform \& |[active]| enthalten \& |[active]| enthalten \& |[active]| enthalten \& |[active]| enthalten \& |[optional]| optional \\
    };
  \end{tikzpicture}}
  \caption{Strukturelle Feature-Zuordnung der Projektprofile}
  \label{fig:profile-matrix}
\end{figure}

\begin{table}[H]
  \centering
  \small
  \renewcommand{\arraystretch}{1.2}
  \caption{Struktur zur Dokumentation der Projektprofile}
  \label{tab:profile-matrix}
  \begin{tabularx}{\textwidth}{@{}p{0.19\textwidth}p{0.19\textwidth}p{0.27\textwidth}>{\raggedright\arraybackslash}X@{}}
    \toprule
    \textbf{Profil} & \textbf{Plattformbezug} & \textbf{enthaltene Komponenten} & \textbf{optionale Capabilities} \\
    \midrule
    % Später Profildefinitionen aus den Projektartefakten dokumentieren.
    % PostgreSQL nur bei backendfähigen Profilen als optionale Capability führen.
    \bottomrule
  \end{tabularx}
\end{table}

